\setcounter{chapter}{10}
\chapter{Кейс-галерея: разбор известных персонажей как комбинаций осей}

\section{Rei Ayanami: kuudere, doll, artificial girl, trauma, silence}

Rei уже была нашим главным примером нулевой точки (1.4), пустых глаз (6.3), экономной речи (7.3) и empty doll предостережения (9.3) — здесь стоит не повторять эти наблюдения, а собрать их в единый вектор и добавить то, что ещё не звучало.

Вектор: affect — минимальный, почти нулевой по умолчанию; роль — сначала инструмент отцовского проекта, затем, к финалу, что-то похожее на дочь и союзницу; intimacy — избегающая, но не защитно, как kuudere (3.4), а буквально, потому что вектор близости у неё изначально не определён; agency — от нуля до слабого, но растущего "партнёр"; conflict — почти отсутствует в начале, появляется как способность отказывать ближе к концу истории.

Деталь, которая обычно теряется за разговором про её "пустоту": Rei не начинает с нуля произвольно — она клон, одна из нескольких идентичных копий, и её путь к индивидуальности проходит буквально через выбор быть именно этой копией, а не взаимозаменяемой единицей. Это придаёт её нулевой точке конкретную, а не абстрактную ставку: обретение вектора для неё равно обретению права быть незаменимой.

Для книги Rei — эталонный кейс того, что 9.3 требует "труда, оправдывающего пустоту": путь от нуля до направления у неё не случаен и не быстр, он занимает практически весь сериал, и именно поэтому работает как модель, а не как трюк.
